%%======================================================================
%%  NT-4a: Literature Extraction & Convention Fixing
%%         for Linearized Nonlocal Field Equations
%%
%%  This document collects formulas from the literature needed to
%%  derive the linearized gravitational field equations from the
%%  SCT spectral action, including propagator structure, Barnes-Rivers
%%  projectors, and the modified Newtonian potential.
%%
%%  RESEARCH DOCUMENT: formulas extracted from cited papers.
%%  No original derivations --- gaps noted explicitly.
%%======================================================================
\documentclass[12pt,a4paper]{article}

%% ---- packages -------------------------------------------------------
\usepackage[utf8]{inputenc}
\usepackage[T1]{fontenc}
\usepackage{lmodern}
\usepackage{amsmath,amssymb,amsthm,mathtools}
\usepackage[margin=2.5cm]{geometry}
\usepackage{booktabs}
\usepackage{enumitem}
\usepackage{hyperref}
\usepackage{xcolor}
\usepackage{fancyhdr}
\usepackage{titlesec}
\usepackage{longtable}

%% ---- theorem-like environments --------------------------------------
\newtheorem{proposition}{Proposition}[section]
\newtheorem{lemma}[proposition]{Lemma}
\newtheorem{corollary}[proposition]{Corollary}
\newtheorem{theorem}[proposition]{Theorem}
\theoremstyle{definition}
\newtheorem{definition}[proposition]{Definition}
\newtheorem{convention}[proposition]{Convention}
\theoremstyle{remark}
\newtheorem{remark}[proposition]{Remark}

%% ---- custom commands ------------------------------------------------
\newcommand{\Tr}{\operatorname{Tr}}
\newcommand{\tr}{\operatorname{tr}}
\newcommand{\Rsq}{R_{\mu\nu\rho\sigma}R^{\mu\nu\rho\sigma}}
\newcommand{\Ricsq}{R_{\mu\nu}R^{\mu\nu}}
\newcommand{\Weylsq}{C_{\mu\nu\rho\sigma}C^{\mu\nu\rho\sigma}}
\newcommand{\dd}{\mathrm{d}}
\newcommand{\ii}{\mathrm{i}}
\newcommand{\Id}{\mathrm{Id}}
\newcommand{\erfi}{\operatorname{erfi}}
\DeclareMathOperator{\sgn}{sgn}

%% ---- annotation commands --------------------------------------------
\newcommand{\fromlit}[1]{\textcolor{blue!70!black}{\textsf{[#1]}}}
\newcommand{\oursign}{\textcolor{teal!80!black}{\textsf{[our convention]}}}
\newcommand{\gap}{\textcolor{red!80!black}{\textsf{[GAP]}}}
\newcommand{\needscheck}{\textcolor{orange!80!black}{\textsf{[needs check]}}}
\newcommand{\exactlit}{\textcolor{blue!60!black}{\textsf{[exact from lit.]}}}
\newcommand{\verified}[1]{%
  \par\smallskip\noindent
  \colorbox{green!10}{\parbox{\dimexpr\linewidth-2\fboxsep}{%
    \textcolor{green!50!black}{\textbf{[Verified]}} #1}}%
  \par\smallskip}
\newcommand{\signcorrection}[1]{%
  \par\noindent\colorbox{red!10}{%
    \parbox{\dimexpr\linewidth-2\fboxsep}{%
      \textcolor{red!50!black}{\textbf{[Sign correction]}}
      \textcolor{red!50!black}{#1}%
    }%
  }\par%
}

%% ---- page style -----------------------------------------------------
\pagestyle{fancy}
\fancyhf{}
\fancyhead[L]{\small\textsc{NT-4a: Literature \& Conventions for Linearized Field Equations}}
\fancyhead[R]{\small\thepage}
\renewcommand{\headrulewidth}{0.4pt}

%% ---- title ----------------------------------------------------------
\title{\textbf{NT-4a: Literature Extraction \& Convention Fixing\\[4pt]
       for Linearized Nonlocal Field Equations}}
\author{David Alfyorov}
\date{March 2026 --- Working Document}

%%======================================================================
\begin{document}
\maketitle

\begin{center}
\small\textit{Credits: Aliaksandr Samatyia contributed research-assistance and workflow support.}
\end{center}

\thispagestyle{empty}

\begin{abstract}
This document collects all formulas, conventions, and computational
recipes from the literature needed to derive the linearized
gravitational field equations from the SCT spectral action
$\Tr(f(D^2/\Lambda^2))$.  The linearization is performed on a flat
Minkowski (or Euclidean) background, yielding the modified graviton
propagator in momentum space.  The document covers: the structure of
higher-derivative gravity (Stelle 1977), the Barnes--Rivers spin
projector decomposition (van Nieuwenhuizen 1973), nonlocal propagator
modifications (BMS, Modesto), and the modified Newtonian potential.
All formulas are extracted from the literature; no original derivations
are performed.  Each formula is annotated with its source.
Gaps are noted explicitly.

\medskip\noindent
\textbf{Key references:}
\begin{itemize}[nosep]
\item Stelle: Phys.\ Rev.\ D \textbf{16} (1977) 953.
\item van Nieuwenhuizen: Nucl.\ Phys.\ B \textbf{60} (1973) 478.
\item Barvinsky--Vilkovisky (BV): Nucl.\ Phys.\ B \textbf{282} (1987)
  163; B \textbf{333} (1990) 471.
\item Biswas--Mazumdar--Siegel (BMS): hep-th/0508194.
\item Modesto: arXiv:1107.2403 [hep-th].
\item Biswas--Gerber--Mazumdar (BGM): arXiv:1112.5430 [gr-qc].
\item Codello--Zanusso (CZ): arXiv:1203.2034 [hep-th].
\item Codello--Percacci--Rahmede (CPR): arXiv:0805.2909 [hep-th].
\item Chamseddine--Connes (CC): hep-th/9606001.
\end{itemize}
\end{abstract}

\tableofcontents
\newpage

%%======================================================================
\section{Background: From NT-1b Phase~3 to NT-4a}
\label{sec:background}
%%======================================================================

\subsection{Input from preceding derivations}

The NT-4a linearization takes as input the one-loop effective action
established in NT-1 through NT-1b Phase~3:
\begin{equation}
  \Gamma_{\text{1-loop}} = \int\dd^4x\,\sqrt{g}\left[
    F_1\!\left(\frac{\Box}{\Lambda^2}\right)\Weylsq
    + F_2\!\left(\frac{\Box}{\Lambda^2}\right)R^2
  \right] + \text{topological},
  \label{eq:effective-action}
\end{equation}
where the form factors $F_1(z)$ and $F_2(z)$ are entire functions
(NT-2) with known local coefficients \fromlit{NT-1b Phase~3}:
\begin{align}
  F_1(0) &= \frac{\alpha_C}{16\pi^2}
  = \frac{13}{1920\pi^2}\,,
    \label{eq:F1-zero}\\
  F_2(0;\xi) &= \frac{\alpha_R(\xi)}{16\pi^2}
  = \frac{2(\xi-1/6)^2}{16\pi^2}\,.
    \label{eq:F2-zero}
\end{align}

\begin{remark}[Curvature basis]
\label{rem:curvature-basis}
\fromlit{CZ Eq.~(25); NT-1b Phase~3}
The Weyl-$R^2$ basis is preferred over the Ricci-$R^2$ basis because:
\begin{enumerate}[nosep]
  \item $\Weylsq$ and $R^2$ are algebraically independent;
  \item The Weyl tensor is traceless ($C^\mu{}_{\mu\rho\sigma}=0$),
    so $\Weylsq$ projects purely onto spin-2;
  \item In the Gauss--Bonnet identity
    $E_4 = \Weylsq - 2\Ricsq + \frac{2}{3}R^2$, the Euler density
    $E_4$ is topological in 4D and does not contribute to
    equations of motion.
\end{enumerate}
The conversion between bases is:
\begin{equation}
  \Ricsq = \Weylsq - \frac{1}{3}R^2 + E_4\,,
  \qquad\text{(in 4D)}.
  \label{eq:basis-conversion}
\end{equation}
\end{remark}


\subsection{Connection to the Stelle framework}

\begin{convention}[Local limit coefficients]
\label{conv:local-coeffs}
\oursign\;
The local ($z\to 0$) limit of the effective action
\eqref{eq:effective-action} is:
\begin{equation}
  \Gamma_{\text{local}} = \int\dd^4x\,\sqrt{g}\left[
    \frac{\alpha_C}{16\pi^2}\,\Weylsq
    + \frac{\alpha_R}{16\pi^2}\,R^2
  \right].
  \label{eq:local-limit}
\end{equation}
Comparing with the Stelle action
$S_{\text{Stelle}} = \int(c_1 R^2 + c_2 \Ricsq)$
and using \eqref{eq:basis-conversion}, we identify
\fromlit{Stelle 1977; NT-1b Phase~3}:
\begin{align}
  c_2 &= 2\alpha_C = \frac{13}{60}\,,
    \label{eq:c2-local}\\
  c_1 &= \frac{\alpha_R - \alpha_C/3}{16\pi^2}\,,\notag\\
  3c_1 + c_2 &= 6\left(\xi - \frac{1}{6}\right)^2.
    \label{eq:3c1c2-local}
\end{align}
\end{convention}

\verified{$c_2 = 13/60$, $3c_1+c_2 = 6(\xi-1/6)^2$ confirmed by
NT-1b Phase~3 (354/354 checks PASS).  Lean~4 formal proof:
\texttt{Propagator.lean}, Theorem \texttt{c2\_local}.}


%%======================================================================
\section{Linearization: Conventions and Identities}
\label{sec:linearization}
%%======================================================================

\subsection{Metric perturbation}

\begin{convention}[Background plus perturbation]
\label{conv:metric-perturbation}
\oursign\;
We expand the metric about a flat background:
\begin{equation}
  g_{\mu\nu} = \eta_{\mu\nu} + h_{\mu\nu}\,,
  \qquad |h_{\mu\nu}| \ll 1\,,
  \label{eq:metric-expansion}
\end{equation}
with $\eta_{\mu\nu} = \text{diag}(-1,+1,+1,+1)$ in Lorentzian
signature or $\delta_{\mu\nu}$ in Euclidean.  We work to
\textbf{first order in $h_{\mu\nu}$} (linear perturbation theory).
\end{convention}

\begin{convention}[Transverse-traceless gauge]
\label{conv:tt-gauge}
\oursign\;
The TT gauge conditions:
\begin{equation}
  \partial^\mu h_{\mu\nu} = 0\,,
  \qquad h^\mu{}_\mu = 0\,.
  \label{eq:tt-conditions}
\end{equation}
This removes the gauge and trace degrees of freedom, leaving
only the physical spin-2 graviton modes.
\end{convention}


\subsection{Linearized curvature tensors}

\begin{proposition}[Linearized curvatures on flat background]
\label{prop:linearized-curvatures}
\fromlit{Wald Ch.~7; Stelle 1977 App.}
To first order in $h_{\mu\nu}$ about $\eta_{\mu\nu}$:
\begin{align}
  R_{\mu\nu}^{(1)} &= \frac{1}{2}\bigl(
    -\Box h_{\mu\nu} + \partial_\mu\partial_\rho h^\rho{}_\nu
    + \partial_\nu\partial_\rho h^\rho{}_\mu
    - \partial_\mu\partial_\nu h
  \bigr),
  \label{eq:ricci-lin}\\[4pt]
  R^{(1)} &= \partial_\mu\partial_\nu h^{\mu\nu} - \Box h\,,
  \label{eq:scalar-lin}\\[4pt]
  G_{\mu\nu}^{(1)} &= R_{\mu\nu}^{(1)}
    - \frac{1}{2}\eta_{\mu\nu}R^{(1)}\,,
  \label{eq:einstein-lin}
\end{align}
where $h = \eta^{\mu\nu}h_{\mu\nu}$ and $\Box =
\eta^{\mu\nu}\partial_\mu\partial_\nu$.
\end{proposition}

\begin{corollary}[TT gauge simplification]
\label{cor:tt-simplification}
In TT gauge ($\partial^\mu h_{\mu\nu} = 0$, $h = 0$):
\begin{align}
  R_{\mu\nu}^{(1)}\big|_{\text{TT}} &= -\frac{1}{2}\Box h_{\mu\nu}\,,
  \label{eq:ricci-tt}\\
  R^{(1)}\big|_{\text{TT}} &= 0\,,
  \label{eq:scalar-tt}\\
  G_{\mu\nu}^{(1)}\big|_{\text{TT}} &= -\frac{1}{2}\Box h_{\mu\nu}\,.
  \label{eq:einstein-tt}
\end{align}
In particular, $R^2|_{\text{TT}} = 0$, so the $F_2(\Box/\Lambda^2)R^2$
term does not contribute in TT gauge.  Only the Weyl term survives.
\end{corollary}

\begin{remark}[Curvature-squared in TT gauge]
In TT gauge, the linearized Weyl tensor coincides with the Riemann
tensor, and \fromlit{Stelle App.~A}:
\begin{equation}
  \Weylsq\big|_{\text{TT}}^{(2)}
  = \frac{1}{2}\,h_{\mu\nu}\,\Box^2\, h^{\mu\nu}
  + \text{total derivatives}\,.
  \label{eq:weyl-sq-tt}
\end{equation}
This gives the fourth-derivative kinetic term characteristic of
higher-derivative gravity.
\end{remark}


\subsection{Off-shell gauge invariance}

\begin{definition}[Linearized gauge transformation]
\label{def:gauge-transf}
\fromlit{Wald Ch.~7; Weinberg Ch.~10}
Under an infinitesimal diffeomorphism generated by $\xi^\mu$:
\begin{equation}
  h_{\mu\nu} \to h_{\mu\nu} + \partial_\mu\xi_\nu
  + \partial_\nu\xi_\mu\,.
  \label{eq:gauge-transf}
\end{equation}
A gauge mode is $h_{\mu\nu}^{(\xi)} = \partial_\mu\xi_\nu
+ \partial_\nu\xi_\mu$.
\end{definition}

\begin{proposition}[Linearized Bianchi identity]
\label{prop:bianchi}
\exactlit\;
The linearized Einstein tensor satisfies
\begin{equation}
  \partial^\mu G_{\mu\nu}^{(1)} = 0
  \label{eq:bianchi}
\end{equation}
identically (not just on-shell).  This is the linearized version
of $\nabla^\mu G_{\mu\nu} = 0$.
\end{proposition}
\begin{proof}
Direct computation from \eqref{eq:ricci-lin}--\eqref{eq:einstein-lin}:
\begin{equation}
  \partial^\mu G_{\mu\nu}^{(1)}
  = \frac{1}{2}\bigl(
    -\Box\partial_\rho h^\rho{}_\nu
    + \partial_\nu\partial_\mu\partial_\rho h^{\mu\rho}
    + \Box\partial_\rho h^\rho{}_\nu
    - \partial_\nu\Box h
    - \partial_\nu\partial_\mu\partial_\rho h^{\mu\rho}
    + \partial_\nu\Box h
  \bigr) = 0\,.
\end{equation}
All six terms cancel pairwise.
\end{proof}

\begin{corollary}[Off-shell gauge invariance]
\label{cor:gauge-inv}
The linearized Einstein tensor is invariant under
\eqref{eq:gauge-transf}:
\begin{equation}
  G_{\mu\nu}^{(1)}[h + \partial\xi + \partial\xi]
  = G_{\mu\nu}^{(1)}[h]\,,
\end{equation}
because $G_{\mu\nu}^{(1)}[\partial\xi + \partial\xi] = 0$ by the
Bianchi identity and the symmetry of the construction.
\end{corollary}

\verified{Off-shell gauge invariance verified numerically for 5 random
$k$-vectors (seed 42) and 5 random $\xi$-vectors: all residuals
$< 10^{-14}$.  Bianchi identity verified for 5 random $k$-vectors
and 5 random symmetric $h$-tensors (seed 43).}


%%======================================================================
\section{Barnes--Rivers Projectors}
\label{sec:barnes-rivers}
%%======================================================================

\subsection{Spin decomposition of the graviton}

\begin{definition}[Symmetric tensor decomposition]
\label{def:spin-decomposition}
\fromlit{van Nieuwenhuizen, Nucl.\ Phys.\ B \textbf{60} (1973) 478;
Rivers, Nuovo Cimento \textbf{34} (1964) 387}
A symmetric rank-2 tensor $h_{\mu\nu}$ in $d=4$ momentum space
decomposes into irreducible spin components:
\begin{itemize}[nosep]
  \item \textbf{Spin-2:} 5 d.o.f.\ (transverse-traceless)
  \item \textbf{Spin-1:} 4 d.o.f.\ (gauge, pure diffeomorphism)
  \item \textbf{Spin-0 (scalar):} 1 d.o.f.\ (trace + longitudinal)
\end{itemize}
Total: $10 = 5 + 4 + 1$ components.
\end{definition}


\subsection{The five projectors}

\begin{convention}[Transverse projector]
\label{conv:theta}
\oursign\;
\fromlit{van Nieuwenhuizen Eq.~(2.1)}
\begin{equation}
  \theta_{\mu\nu}(k) = \delta_{\mu\nu}
  - \frac{k_\mu k_\nu}{k^2}\,.
  \label{eq:theta}
\end{equation}
In Euclidean signature, $\theta$ projects onto the subspace
transverse to $k^\mu$.  We have $\theta_\mu{}^\mu = d-1 = 3$ and
$k^\mu\theta_{\mu\nu} = 0$.
\end{convention}

\begin{definition}[Barnes--Rivers projectors]
\label{def:barnes-rivers}
\fromlit{van Nieuwenhuizen Eqs.~(2.3)--(2.7)}
The five orthogonal projectors on symmetric rank-2 tensors
in momentum space are:
\begin{align}
  P^{(2)}_{\mu\nu,\rho\sigma}
  &= \frac{1}{2}\bigl(
    \theta_{\mu\rho}\theta_{\nu\sigma}
    + \theta_{\mu\sigma}\theta_{\nu\rho}
  \bigr) - \frac{1}{3}\theta_{\mu\nu}\theta_{\rho\sigma}\,,
    \label{eq:P2}\\[4pt]
  P^{(1)}_{\mu\nu,\rho\sigma}
  &= \frac{1}{2}\bigl(
    \theta_{\mu\rho}\omega_{\nu\sigma}
    + \theta_{\mu\sigma}\omega_{\nu\rho}
    + \theta_{\nu\rho}\omega_{\mu\sigma}
    + \theta_{\nu\sigma}\omega_{\mu\rho}
  \bigr),
    \label{eq:P1}\\[4pt]
  P^{(0\text{-}s)}_{\mu\nu,\rho\sigma}
  &= \frac{1}{3}\theta_{\mu\nu}\theta_{\rho\sigma}\,,
    \label{eq:P0s}\\[4pt]
  P^{(0\text{-}w)}_{\mu\nu,\rho\sigma}
  &= \omega_{\mu\nu}\omega_{\rho\sigma}\,,
    \label{eq:P0w}\\[4pt]
  P^{(0\text{-}sw)}_{\mu\nu,\rho\sigma}
  &= \frac{1}{\sqrt{3}}\bigl(
    \theta_{\mu\nu}\omega_{\rho\sigma}
    + \omega_{\mu\nu}\theta_{\rho\sigma}
  \bigr),
    \label{eq:P0sw}
\end{align}
where $\omega_{\mu\nu} = k_\mu k_\nu/k^2$ is the longitudinal
projector ($\theta + \omega = \delta$).
\end{definition}

\begin{proposition}[Projector algebra]
\label{prop:projector-algebra}
\exactlit\;
\fromlit{van Nieuwenhuizen Eqs.~(2.8)--(2.10)}
The Barnes--Rivers projectors satisfy:
\begin{enumerate}[nosep]
  \item \textbf{Orthogonality:}
    $P^{(A)}_{\mu\nu,\alpha\beta}\,P^{(B)\alpha\beta}{}_{\rho\sigma}
    = \delta^{AB}\,P^{(A)}_{\mu\nu,\rho\sigma}$\,.
  \item \textbf{Completeness:}
    $P^{(2)} + P^{(1)} + P^{(0\text{-}s)} + P^{(0\text{-}w)}
    = I_{\mu\nu,\rho\sigma}$\,,
    where $I_{\mu\nu,\rho\sigma} = \frac{1}{2}(\delta_{\mu\rho}
    \delta_{\nu\sigma} + \delta_{\mu\sigma}\delta_{\nu\rho})$.
  \item \textbf{Traces:}
    $P^{(2)\mu\nu}{}_{\mu\nu} = 5$,
    $P^{(1)\mu\nu}{}_{\mu\nu} = 4$,
    $P^{(0\text{-}s)\mu\nu}{}_{\mu\nu} = 1$.
\end{enumerate}
Note: $P^{(0\text{-}sw)}$ is \emph{not} a projector (it is
nilpotent: $P^{(0\text{-}sw)}\cdot P^{(0\text{-}sw)} = 0$).
It is the ``mixing'' operator between the two spin-0 sectors.
\end{proposition}

\verified{Projector algebra (orthogonality, completeness, traces)
verified numerically for 5 random $k$-vectors at double precision.
Residuals $< 10^{-14}$.  Lean~4: \texttt{GaugeInv.lean},
Theorems \texttt{projector\_P2\_squared},
\texttt{projector\_P2\_P0s\_orthogonal},
\texttt{P2\_plus\_P0s\_trace}.}


\subsection{Role in the graviton propagator}

\begin{remark}[Why Barnes--Rivers projectors matter]
\fromlit{van Nieuwenhuizen 1973; Stelle 1977}
The graviton propagator in any higher-derivative or nonlocal theory
can be decomposed as:
\begin{equation}
  G_{\mu\nu,\rho\sigma}(k) = \frac{P^{(2)}_{\mu\nu,\rho\sigma}}{
    k^2\,\Pi_{\text{TT}}(k^2/\Lambda^2)}
  + \frac{P^{(0\text{-}s)}_{\mu\nu,\rho\sigma}}{
    k^2\,\Pi_s(k^2/\Lambda^2)}
  + \text{gauge-dependent terms}\,.
  \label{eq:propagator-decomposition}
\end{equation}
The physical content is entirely in the spin-2 and spin-0-scalar
denominators $\Pi_{\text{TT}}$ and $\Pi_s$.  The gauge-dependent
terms (involving $P^{(1)}$, $P^{(0\text{-}w)}$, $P^{(0\text{-}sw)}$)
are unphysical and drop out of on-shell amplitudes.
\end{remark}


%%======================================================================
\section{Stelle's Local Theory: Propagator Structure}
\label{sec:stelle}
%%======================================================================

\subsection{The quadratic-gravity action}

\begin{convention}[Stelle action]
\label{conv:stelle-action}
\fromlit{Stelle, Phys.\ Rev.\ D \textbf{16} (1977) 953, Eq.~(1)}
The most general local action quadratic in curvature (in 4D):
\begin{equation}
  S_{\text{Stelle}} = \int\dd^4x\,\sqrt{g}\left[
    \frac{R}{16\pi G}
    + \alpha\,\Weylsq + \beta\,R^2
  \right].
  \label{eq:stelle-action}
\end{equation}
\end{convention}

\begin{remark}[Relation to our notation]
In the SCT local limit: $\alpha = F_1(0)$ and $\beta = F_2(0)$.
Stelle uses different notation; the mapping is:
\begin{center}
\begin{tabular}{lcc}
\toprule
Quantity & Stelle & SCT \\
\midrule
Weyl$^2$ coefficient & $\alpha$ & $F_1(0) = \alpha_C/(16\pi^2)$ \\
$R^2$ coefficient & $\beta$ & $F_2(0) = \alpha_R/(16\pi^2)$ \\
\bottomrule
\end{tabular}
\end{center}
\end{remark}


\subsection{Stelle propagator}

\begin{proposition}[Stelle propagator in Barnes--Rivers basis]
\label{prop:stelle-propagator}
\fromlit{Stelle 1977, Eq.~(15)}
Linearizing $S_{\text{Stelle}}$ about flat space and inverting
the kinetic operator:
\begin{align}
  G^{\text{Stelle}}_{\text{TT}}(k^2)
  &= \frac{1}{k^2(1 + 2\alpha\, k^2)}
  = \frac{1}{k^2} - \frac{1}{k^2 + m_2^2}\,,
  \label{eq:stelle-prop-tt}\\[4pt]
  G^{\text{Stelle}}_s(k^2)
  &= \frac{1}{k^2(1 + (6\beta + 2\alpha)\, k^2)}\,,
  \label{eq:stelle-prop-s}
\end{align}
where $m_2^2 = 1/(2\alpha)$ is the massive spin-2 ghost mass.
\end{proposition}

\begin{remark}[The ghost]
\fromlit{Stelle 1977, \S III}
The partial fraction decomposition \eqref{eq:stelle-prop-tt}
shows a massless graviton pole at $k^2=0$ (positive residue)
and a massive pole at $k^2 = -m_2^2$ (\textbf{negative} residue).
The negative residue means the massive spin-2 mode is a ghost:
it has negative-norm states and violates unitarity.

This is the central problem that the nonlocal structure of the
spectral action aims to resolve: by replacing $2\alpha \to
c_2\, F_1(k^2/\Lambda^2)/F_1(0)$, the denominator
$\Pi_{\text{TT}}(z) = 1 + c_2\, z\, \hat{F}_1(z)$ can be made
zero-free if $\hat{F}_1$ is entire and positive on the real axis.
\end{remark}

\signcorrection{The factor $2\alpha$ in the Stelle propagator
denominator maps to $c_2 = 2\alpha_C$ in our notation.
Some references write $c_2 = \alpha_C$; the factor of 2 comes
from the conversion between the Weyl-$R^2$ basis and the
Ricci-$R^2$ basis via $\Ricsq = \Weylsq - R^2/3 + E_4$.}


%%======================================================================
\section{Nonlocal Propagator: SCT Structure}
\label{sec:nonlocal-propagator}
%%======================================================================

\subsection{Modified denominators}

\begin{definition}[SCT propagator denominators]
\label{def:sct-denominators}
\oursign\;
\fromlit{NT-4a derivation; cf.\ BMS Eq.~(7)}
Define the normalized shape functions:
\begin{equation}
  \hat{F}_i(z) = \frac{F_i(z)}{F_i(0)}\,,
  \qquad \hat{F}_i(0) = 1\,.
  \label{eq:shape-functions}
\end{equation}
The modified propagator denominators are:
\begin{align}
  \Pi_{\text{TT}}(z) &= 1 + c_2\, z\, \hat{F}_1(z)
  = 1 + \frac{13}{60}\, z\, \hat{F}_1(z)\,,
    \label{eq:Pi-TT}\\
  \Pi_s(z;\xi) &= 1 + (3c_1+c_2)\, z\, \hat{F}_2(z;\xi)
  = 1 + 6\left(\xi-\frac{1}{6}\right)^2 z\, \hat{F}_2(z;\xi)\,.
    \label{eq:Pi-s}
\end{align}
\end{definition}

\begin{remark}[Normalization convention]
\label{rem:normalization}
The normalization $\hat{F}_i(0) = 1$ ensures
$\Pi_{\mathrm{TT}}(0) = \Pi_s(0) = 1$, reproducing
the standard GR propagator at zero momentum.
This is the same convention used by BMS~\cite{BMS2006}
and Modesto~\cite{Modesto2012}.
\end{remark}

\verified{$\Pi_{\text{TT}}(0) = 1$ and $\Pi_s(0) = 1$ confirmed
analytically and numerically (100-digit precision).
Lean~4: \texttt{Propagator.lean}, Theorem \texttt{pi\_tt\_zero}.}


\subsection{Comparison with BMS and Modesto}

\begin{remark}[BMS ghost-free ansatz]
\fromlit{BMS hep-th/0508194, Eq.~(7)}
BMS propose the form $a(-\Box/M^2) = e^{-\Box/M^2}$, giving
$\Pi = 1 + c\, k^2\, e^{k^2/M^2}$.  This is manifestly positive
for all $k^2 \geq 0$ and has no zeros anywhere.

In contrast, the SCT form factors are \emph{not} pure exponentials:
$\hat{F}_1(z) \neq e^{z/M^2}$.  The SCT form factors are
\emph{derived} from the spectral action, not chosen by hand.
Ghost-freedom must therefore be verified rather than assumed.
\end{remark}

\begin{remark}[Modesto super-renormalizability]
\fromlit{Modesto 1107.2403}
Modesto showed that choosing $a(\Box) = e^{H(\Box)}$ with $H$
a polynomial of sufficiently high degree leads to
super-renormalizable (or even finite) quantum gravity.  The SCT
spectral action produces a similar UV improvement, but through
the heat kernel mechanism rather than by hand.
\end{remark}


%%======================================================================
\section{Scalar Mode and the Higgs Non-Minimal Coupling}
\label{sec:scalar-mode}
%%======================================================================

\begin{proposition}[Scalar mode decoupling at $\xi = 1/6$]
\label{prop:scalar-decoupling}
\fromlit{NT-1b Phase~1; Stelle 1977 \S III}
At conformal coupling $\xi = 1/6$:
\begin{align}
  \alpha_R(1/6) &= 0\,,\\
  3c_1 + c_2\big|_{\xi=1/6} &= 0\,,\\
  \Pi_s(z;\xi=1/6) &= 1 \quad\text{(identically)}\,.
\end{align}
The scalar gravitational mode has the same propagator as in GR
(no higher-derivative modification).  The scalar mode is
\textbf{physical but unmodified} --- it propagates at all momenta
without a mass scale.
\end{proposition}

\begin{remark}[Conformal vs.\ minimal coupling]
The two special values of $\xi$ are:
\begin{center}
\begin{tabular}{lcc}
\toprule
Coupling & $\xi$ & $3c_1+c_2$ \\
\midrule
Conformal & $1/6$ & $0$ (scalar decouples) \\
Minimal & $0$ & $1/6$ (scalar mass $m_0 = \Lambda\sqrt{6}$) \\
\bottomrule
\end{tabular}
\end{center}
At minimal coupling, the scalar mode acquires a mass
$m_0^2 = \Lambda^2/(3c_1+c_2) = 6\Lambda^2$.  Whether this
mode is ghost or healthy depends on the sign of the residue.
\end{remark}

\verified{$\Pi_s(z;\xi=1/6) = 1$ verified analytically and
numerically for $z \in \{0.1, 0.5, 1, 5, 10\}$.
Lean~4: \texttt{Propagator.lean}, Theorems
\texttt{scalar\_mode\_conformal} and \texttt{scalar\_mode\_minimal}.}


%%======================================================================
\section{Modified Newtonian Potential}
\label{sec:newtonian}
%%======================================================================

\subsection{Fourier representation}

\begin{convention}[Static potential from propagator]
\label{conv:potential}
\fromlit{Stelle 1977 Eq.~(17); BGM 1112.5430 Eq.~(12)}
The gravitational potential of a static point source of mass $M$
in the linearized theory is obtained by Fourier transform:
\begin{equation}
  V(r) = -\frac{G M}{r}\int_0^\infty\dd k\,
  \frac{k\sin(kr)}{\pi}\left[
    \frac{4}{3}\,G_{\text{TT}}(k^2)
    - \frac{1}{3}\,G_s(k^2)
  \right]
  \cdot k^2\,.
  \label{eq:potential-fourier}
\end{equation}
The factors $4/3$ and $-1/3$ come from the spin-2 and spin-0
contributions to the Newtonian limit, respectively.
\end{convention}

\begin{remark}[Alternative form]
\fromlit{BGM 1112.5430}
For a nonlocal propagator $G(k^2) = 1/(k^2\,\Pi(k^2/\Lambda^2))$,
the potential ratio $V(r)/V_{\text{Newton}}(r)$ can be written as
\begin{equation}
  \frac{V(r)}{V_N(r)} = 1 + \Delta V_{\text{TT}}(r)
  + \Delta V_s(r)\,,
  \label{eq:potential-ratio}
\end{equation}
where $\Delta V_{\text{TT}}$ and $\Delta V_s$ are the spin-2 and
spin-0 corrections, respectively.
\end{remark}


\subsection{Yukawa approximation}

\begin{proposition}[Effective Yukawa form]
\label{prop:yukawa}
\fromlit{Stelle 1977 Eq.~(18); BGM Eq.~(14)}
When the propagator denominators have effective mass scales
$m_2$ and $m_0$, the leading correction to the Newtonian potential
takes the Yukawa form:
\begin{equation}
  \frac{V(r)}{V_N(r)} \approx
  1 - \frac{4}{3}\,e^{-m_2 r}
  + \frac{1}{3}\,e^{-m_0 r}\,,
  \label{eq:yukawa-approx}
\end{equation}
where the signs reflect the ghost nature of the Stelle spin-2 mode.
At $r \gg 1/m_2, 1/m_0$, both Yukawa terms are exponentially
suppressed and $V(r) \to V_N(r)$.
\end{proposition}

\begin{remark}[SCT mass scales]
\fromlit{NT-4a derivation}
The effective mass scales in the SCT case are:
\begin{align}
  m_2 &= \frac{\Lambda}{\sqrt{c_2}}
  = \Lambda\sqrt{\frac{60}{13}}\,,
    \label{eq:m2}\\
  m_0 &= \frac{\Lambda}{\sqrt{3c_1+c_2}}
  = \frac{\Lambda}{\sqrt{6(\xi-1/6)^2}}\,.
    \label{eq:m0}
\end{align}
At conformal coupling $\xi = 1/6$, $m_0 \to \infty$ (the scalar
Yukawa term vanishes) and
\begin{equation}
  \frac{V(r)}{V_N(r)}\bigg|_{\xi=1/6}
  \approx 1 - \frac{4}{3}\,e^{-m_2 r}\,.
  \label{eq:yukawa-conformal}
\end{equation}
\end{remark}


\subsection{UV behavior of the potential}

\begin{proposition}[Finiteness at the origin]
\label{prop:finite-origin}
\fromlit{Modesto 1107.2403 \S IV; BGM 1112.5430}
Unlike the Newtonian $1/r$ singularity, the modified potential
$V(r)$ is \textbf{finite} as $r \to 0$:
\begin{equation}
  V(0) = -G M \left(\frac{4}{3}\,m_2 - \frac{1}{3}\,m_0\right)
  + O(1/\Lambda^2)\,.
  \label{eq:V-origin}
\end{equation}
This UV regularization is a generic feature of nonlocal
(entire-function) modifications of the propagator.
\end{proposition}

\begin{remark}[Physical significance]
The finiteness of the gravitational potential at $r=0$ is
one of the key predictions of the SCT spectral action:
it resolves the Newtonian singularity without introducing
a physical cutoff or lattice.  The regularization scale
is set by $\Lambda$ (the spectral action cutoff).
\end{remark}

\verified{$V(r)$ finiteness at small $r$ verified numerically:
$V(10^{-6}/\Lambda)$ returns a finite value for both
$\xi = 0$ and $\xi = 1/6$.  Large-$r$ recovery:
$V(r)/V_N(r) \to 1$ with $< 1\%$ error at $r = 10^4/\Lambda$.}


%%======================================================================
\section{Conventions Summary}
\label{sec:summary}
%%======================================================================

\begin{center}
\begin{tabular}{lll}
\toprule
Quantity & Convention & Source \\
\midrule
Metric perturbation & $g_{\mu\nu} = \eta_{\mu\nu}+h_{\mu\nu}$
  & Standard \\
TT gauge & $\partial^\mu h_{\mu\nu}=0$, $h=0$
  & Standard \\
Barnes--Rivers projectors & $P^{(2)}, P^{(1)}, P^{(0\text{-}s)},
  P^{(0\text{-}w)}$ & vN 1973 \\
Transverse projector & $\theta_{\mu\nu} = \delta_{\mu\nu} -
  k_\mu k_\nu/k^2$ & vN 1973 \\
Propagator denominators & $\Pi = 1 + c\cdot z\cdot\hat{F}(z)$
  & BMS/NT-4a \\
Local coefficients & $c_2 = 13/60$,
  $3c_1+c_2=6(\xi-1/6)^2$ & NT-1b Phase~3 \\
Yukawa masses & $m_2 = \Lambda\sqrt{60/13}$,
  $m_0 = \Lambda/\sqrt{6(\xi-1/6)^2}$ & NT-4a derived \\
\bottomrule
\end{tabular}
\end{center}


%%======================================================================
\section{Gap Analysis}
\label{sec:gaps}
%%======================================================================

\begin{enumerate}[label=\textbf{G\arabic*.}]
\item \textbf{Full nonlinear equations:}
  NT-4a covers only the linearized (flat background) case.
  The full nonlinear variational equations
  $\delta S/\delta g^{\mu\nu} = 0$ are deferred to NT-4b
  (Barvinsky--Vilkovisky formalism).

\item \textbf{FLRW reduction:}
  Cosmological applications require reducing the field equations
  to the FLRW metric.  This is NT-4c.

\item \textbf{Gauge fixing:}
  The TT gauge is used throughout.  A covariant gauge-fixing
  analysis (de Donder gauge, etc.) would require the full
  $P^{(1)}$ and $P^{(0\text{-}w)}$ sectors.

\item \textbf{Gravitational waves:}
  The linearized equations describe gravitational wave
  propagation with modified dispersion, but the detailed
  analysis of waveform corrections is not performed here.

\item \textbf{PPN parameters:}
  The post-Newtonian limit and PPN parameter extraction
  from the modified potential is deferred to PPN-1.

\item \textbf{Numerical Fourier transform:}
  The Yukawa approximation is used instead of the full numerical
  Fourier integral \eqref{eq:potential-fourier}.  A numerical
  integration would be needed for precision tests against
  solar system data.
\end{enumerate}


%%======================================================================
\begin{thebibliography}{99}
%%======================================================================

\bibitem{Stelle1977}
  K.~S.~Stelle,
  ``Renormalization of higher-derivative quantum gravity,''
  Phys.\ Rev.\ D \textbf{16} (1977) 953.

\bibitem{vN1973}
  P.~van Nieuwenhuizen,
  ``On the renormalization of quantum gravitation without matter,''
  Annals Phys.\ \textbf{104} (1977) 197;
  see also Nucl.\ Phys.\ B \textbf{60} (1973) 478.

\bibitem{BV1987}
  A.~O.~Barvinsky, G.~A.~Vilkovisky,
  ``Beyond the Schwinger-DeWitt technique: Converting loops into
  trees and in-in currents,''
  Nucl.\ Phys.\ B \textbf{282} (1987) 163.

\bibitem{BV1990a}
  A.~O.~Barvinsky, G.~A.~Vilkovisky,
  ``Covariant perturbation theory (II): Second order in the
  curvature,''
  Nucl.\ Phys.\ B \textbf{333} (1990) 471.

\bibitem{BMS2006}
  T.~Biswas, A.~Mazumdar, W.~Siegel,
  ``Bouncing universes in string-inspired gravity,''
  JCAP \textbf{0603} (2006) 009
  [hep-th/0508194].

\bibitem{Modesto2012}
  L.~Modesto,
  ``Super-renormalizable quantum gravity,''
  Phys.\ Rev.\ D \textbf{86} (2012) 044005
  [arXiv:1107.2403 [hep-th]].

\bibitem{BGM2012}
  T.~Biswas, E.~Gerber, A.~Mazumdar,
  ``Towards singularity- and ghost-free theories of gravity,''
  Phys.\ Rev.\ Lett.\ \textbf{108} (2012) 031101
  [arXiv:1112.5430 [gr-qc]].

\bibitem{CZ2012}
  A.~Codello, A.~Zanusso,
  ``On the non-local heat kernel expansion,''
  J.\ Math.\ Phys.\ \textbf{54} (2013) 013513
  [arXiv:1203.2034 [hep-th]].

\bibitem{CPR2009}
  A.~Codello, R.~Percacci, C.~Rahmede,
  ``Investigating the ultraviolet properties of gravity with a
  Wilsonian renormalization group equation,''
  Annals Phys.\ \textbf{324} (2009) 414
  [arXiv:0805.2909 [hep-th]].

\bibitem{CC1997}
  A.~H.~Chamseddine, A.~Connes,
  ``The spectral action principle,''
  Commun.\ Math.\ Phys.\ \textbf{186} (1997) 731
  [hep-th/9606001].

\bibitem{Tomboulis1997}
  E.~T.~Tomboulis,
  ``Super-renormalizable gauge and gravitational theories,''
  hep-th/9702080 (1997).

\end{thebibliography}

\end{document}
